\section{Conclusion}
\label{sec:conclusion}

This report developed, from first principles, a quaternion-based attitude dynamics and control simulation for a 3U CubeSat and used it to close the loop on a specific, physically motivated maneuver: recovering from a post-deployment tumble of $19.7^\circ$/s while simultaneously executing a $120^\circ$ large-angle slew to a target inertial attitude, using a three-axis reaction wheel assembly and a quaternion PD controller.

Three results, taken together, support the claim that the simulation is a correct implementation of the derived model, not merely a plausible-looking one. First, the closed-loop system converged to within $0.11^\circ$ of the target attitude and $0.01^\circ$/s of zero rate in \SI{74.1}{\second} (\cref{sec:results}), consistent in direction and rough magnitude with the settling time predicted by the small-angle gain design, given the additional gyroscopic coupling and brief actuator saturation the linear design explicitly did not account for (\cref{sec:validation_physical}). Second, three independent numerical checks -- quaternion norm conservation, agreement between two structurally different integrators, and monotonic decrease of the Lyapunov function derived in \cref{sec:pd_law} -- all held to within floating-point or negligible-physical-significance tolerances (\cref{tab:results_summary}). Third, and most tellingly, the Lyapunov monotonicity check was not merely a passive confirmation: it is what surfaced a genuine multiplication-order error in the first implementation of the error quaternion (\cref{sec:results_stability}), an error that produced superficially reasonable-looking convergence but violated an exact, unconditional analytical guarantee. Recovering from that error and re-verifying against the same check is, in the author's judgment, as informative a demonstration of understanding the underlying control theory as the final converged result itself.

The scope of this work was deliberately bounded, and that boundary is repeated here rather than left implicit. This report addressed the rigid-body attitude dynamics and control problem in a deterministic, disturbance-free, truth-model setting, with idealized actuation. It did not address attitude determination (sensor modeling and state estimation), environmental disturbance torque rejection, reaction wheel momentum management, or actuator redundancy -- each identified explicitly in \cref{sec:limitations} as a specific, prioritized direction for extension rather than a vague caveat. Of these, disturbance rejection is judged the most consequential gap: a controller that has only been shown to converge in a torque-free environment has not yet been shown to hold a pointing requirement against the persistent, low-level torques a CubeSat actually experiences on orbit, which is the operational condition the ADCS exists to handle.

Within that bounded scope, the project met the objectives set out in \cref{sec:objectives}: the quaternion kinematics and Euler rigid-body dynamics were derived and explained physically, not just stated; a representative CubeSat mass/inertia model and actuator suite were defined with explicitly stated assumptions; a PD attitude controller was derived and proven almost-globally asymptotically stable; and the resulting closed-loop system was implemented, exercised on a large-angle maneuver chosen specifically to be representative of the class of problem the quaternion representation is preferred for, and validated against both numerical and analytical checks. The most direct extension of this work -- and the one the author would prioritize first -- is adding a disturbance torque model and re-evaluating the same controller's steady-state performance against it, since that is the single change most likely to alter the conclusions drawn here.
